\section*{第 \ref{ch:b2:genfun} 章 --- 概率母函数}

\begin{solution}{exo:b2:genfun:1}
公平骰子：$G(t) = \frac16(t + t^2 + \dots + t^6) = \frac t6(1 + t
+ \dots + t^5)$。若两颗公平骰子的点数和
$\{2, \dots, 12\}$ 上服从均匀分布，则
\[
G(t)^2 = \frac{t^2}{36}\,h(t)^2
= \frac{t^2}{11}\sum_{k=0}^{10}t^k ,
\qquad h(t) = 1 + t + \dots + t^5 .
\]
由于 $h$ 是奇次数 $5$ 的实多项式，因此它有实根
（由介值定理；具体地，$h(-1) = 0$），从而
$h^2$ 有实根。但 $\sum_{k=0}^{10}t^k$ 没有实根：当 $t \geq 0$ 时它为正，
而当 $t < 0$ 时它等于
$\frac{t^{11} - 1}{t - 1}$，是两个负数之商。
矛盾 --- 两颗公平骰子的点数和永远不可能均匀分布（
骰子点数和的熟知三角形分布也说明了这一点）。
\end{solution}

\begin{solution}{exo:b2:genfun:2}
\emph{二项分布：}$G(t) = (1 - p + pt)^n$，$G'(t) = np(1 - p +
pt)^{n-1}$，$G''(t) = n(n-1)p^2(1 - p + pt)^{n-2}$，于是
\[
\E(X) = G'(1) = np,
\qquad
V(X) = G''(1) + G'(1) - G'(1)^2
= n(n-1)p^2 + np - n^2p^2 = np(1-p).
\]
\emph{几何分布}（$q = 1 - p$）：$G(t) = \frac{pt}{1 - qt}$，于是
$G'(t) = \frac{p}{(1 - qt)^2}$，$G''(t) = \frac{2pq}{(1 -
qt)^3}$；取 $t = 1$（利用 $1 - q = p$）：
\[
\E(X) = \frac{p}{p^2} = \frac1p,
\qquad
V(X) = \frac{2q}{p^2} + \frac1p - \frac{1}{p^2}
= \frac{2q + p - 1}{p^2}
= \frac{q}{p^2} ,
\]
与\cref{exo:b2:randomvar:1}所得结果一致，且计算更简便。
\end{solution}

\begin{solution}{exo:b2:genfun:3}
不能，即使两颗骰子的配重不同也不行。设 $X, Y$ 为
$\{1, \dots, 6\}$ 上的分布，其和服从均匀分布。则 $G_X(t) = t\,a(t)$，
$G_Y(t) = t\,b(t)$，其中 $a, b$ 为实多项式，次数\emph{至
多} $5$ --- 且两者次数之和须为 $10$（和以正概率达到
$12$），故 $\deg a = \deg b = 5$，两者
均为奇数。同\cref{exo:b2:genfun:1}，
\[
a(t)\,b(t) = \frac{1}{11}\sum_{k=0}^{10}t^k
\]
将迫使左边有实根（每个奇次实多项式都有一个），而右边没有实根。因此
无论两颗独立骰子
--- 配重是否相同 --- 都不可能使点数之和服从均匀分布。
\end{solution}

\begin{solution}{exo:b2:genfun:4}
由\cref{thm:b2:genfun:compound}，取 $G_N(s) = e^{\lambda(s-1)}$
及 $G_X(t) = 1 - p + pt$：
\[
G_S(t) = e^{\lambda(1 - p + pt - 1)} = e^{\lambda p(t - 1)} :
\]
$S \sim \mathcal{P}(\lambda p)$。舍弃个数 $D = N - S$
所计的是以概率 $1 - p$ 被保留的同一批元素，
由同样的计算可得 $D \sim \mathcal{P}(\lambda(1 - p))$。
独立性可直接验证：对 $j, k \in \N$，
\begin{align*}
\P(S = j,\ D = k)
&= \P(N = j + k)\,\binom{j+k}{j}p^jq^k
= e^{-\lambda}\frac{\lambda^{j+k}}{(j+k)!}\,
\frac{(j+k)!}{j!\,k!}\,p^jq^k\\
&= \Bigl(e^{-\lambda p}\frac{(\lambda p)^j}{j!}\Bigr)
\Bigl(e^{-\lambda q}\frac{(\lambda q)^k}{k!}\Bigr)
\end{align*}
其中 $q = 1 - p$：联合分布可分解为
$\mathcal{P}(\lambda p) \otimes \mathcal{P}(\lambda q)$。泊松流
经随机分流得到\emph{独立}的泊松流 ---
这是排队论中一个常用的小小奇迹。
\end{solution}

\begin{solution}{exo:b2:genfun:5}
相邻两次出现正面之间的等待时间是独立的
服从几何分布 $\mathcal{G}(p)$ 的随机变量（无记忆性：每次出现
正面后游戏重新开始），故 $T_r = W_1 + \dots + W_r$，
由乘积性（\cref{thm:b2:genfun:product}）可得
\[
G_{T_r}(t) = \Bigl(\frac{pt}{1 - qt}\Bigr)^{r},
\qquad
\E(T_r) = r\,\E(W_1) = \frac rp,
\qquad
V(T_r) = r\,V(W_1) = \frac{rq}{p^2}
\]
（$q = 1 - p$；方差由独立性相加）。展开：由
广义二项级数
（\cref{ch:b2:powerseries}），$(1 - qt)^{-r} = \sum_{m\geq0}
\binom{m + r - 1}{r - 1}q^mt^m$，故 $t^n$ 在
$p^rt^r(1 - qt)^{-r}$ 中的系数（取 $m = n - r$）为
\[
\P(T_r = n) = \binom{n-1}{r-1}p^r(1-p)^{n-r},
\qquad n \geq r ,
\]
这就是\emph{负二项分布} --- 组合地看：第 $r$ 次
正面恰落在第 $n$ 次抛掷，当且仅当前 $r - 1$ 次正面
占据了前 $n - 1$ 次抛掷中的位置。
\end{solution}

\begin{solution}{exo:b2:genfun:6}
$m = 1\cdot\frac38 + 2\cdot\frac38 + 3\cdot\frac18 = \frac{3 + 6 +
3}{8} = \frac32 > 1$：超临界。母函数为
\[
G(t) = \frac{1 + 3t + 3t^2 + t^3}{8} = \frac{(1 + t)^3}{8} ,
\]
于是不动点满足 $(1 + t)^3 = 8t$，即 $t^3 + 3t^2 - 5t + 1
= 0$。分出必有的根 $t = 1$：
\[
t^3 + 3t^2 - 5t + 1 = (t - 1)\bigl(t^2 + 4t - 1\bigr),
\]
由 $t^2 + 4t - 1 = 0$ 解得 $t = -2 \pm \sqrt5$。位于
$\intco{0}{1}$ 中的根为 $\sqrt5 - 2 \approx 0.236$，由
\cref{thm:b2:genfun:extinction}，
\[
q = \sqrt 5 - 2 .
\]
（一个令人欣喜的检验：后代分布与 $3$ 枚独立
公平硬币的正面个数同分布，$Z_1 \sim \mathcal{B}(3, \frac12)$。）
\end{solution}

\begin{solution}{exo:b2:genfun:7}
\emph{期望。}首先 $\E(Z_n) = m^n$：由
\cref{thm:b2:genfun:compound}，$\E(Z_{n+1}) = \E(Z_n)\,m$，
且 $\E(Z_0) = 1$。族 $\bigl(Z_n(\omega)\P(\{\omega\})
\bigr)_{n, \omega}$ 非负，因此族的富比尼定理
无条件地适用：
\[
\E(Y) = \sum_{n=0}^{\infty}\E(Z_n)
= \sum_{n=0}^\infty m^n = \frac{1}{1 - m} < \infty
\]
（特别地，$Y$ 几乎必然有限，与
亚临界情形下的必然灭绝相一致）。

\emph{函数方程。}将总体按祖先的
子代分解：若祖先有 $Z_1 = k$ 个子代，则
总后代数为 $Y = 1 + Y_1 + \dots + Y_k$，其中 $Y_i$ 是
第 $i$ 个子代世系的总后代数 --- 且各 $Y_i$ 为
与 $Y$ 独立同分布的副本，且与 $Z_1$ 独立（不同世系
使用互不相交且相互独立的繁殖事件）。以
$Z_1$ 为条件，如\cref{thm:b2:genfun:compound}所示：
\[
H(t) = \E\bigl(t^Y\bigr)
= t\sum_{k=0}^\infty \P(Z_1 = k)\,H(t)^k
= t\,G\bigl(H(t)\bigr),
\]
因子 $t$ 计入祖先自身。（对于
分布 $p_0 = 1 - p$，$p_2 = p$（二分支），则
关于 $H$ 的二次方程可显式求解并展开 ---
\cref{ch:b2:powerseries}中的卡塔兰数计数
家族树。）
\end{solution}

\begin{solution}{exo:b2:genfun:8}
设 $R > 1$ 为收敛半径，取定 $t \in \intoo{1}{R}$。则
$\E(t^X) = G(t) < \infty$，由马尔可夫不等式
（\cref{thm:b2:randomvar:markov}）用于非负
随机变量 $t^X$ 取阈值 $t^n$：
\[
\P(X \geq n) = \P\bigl(t^X \geq t^n\bigr)
\leq \frac{G(t)}{t^n} = C\rho^n,
\qquad C = G(t),\quad \rho = \frac1t \in \intoo{0}{1}.
\]
\emph{逆命题。}若 $\P(X \geq n) \leq C\rho^n$，则 $p_n \leq
\P(X \geq n) \leq C\rho^n$，故当 $\abs t < \frac1\rho$ 时
级数 $\sum p_n\abs t^n$ 被收敛的几何
级数 $C\sum(\rho\abs t)^n$ 所控制，收敛半径至少为
$\frac1\rho > 1$。生成函数的半径与尾部的几何衰减
是同一性质的两个侧面。
\end{solution}

\begin{solution}{exo:b2:genfun:9}
记 $p_k^{(n)} = \P(X_n = k)$，$p_k = \P(X = k)$。

\emph{情形 $k = 0$。}对 $t \in \intoo{0}{1}$ 及任一分布
$(q_j)$ 满足 $\sum_j q_j \leq 1$：
\[
\Bigl|\,q_0 - \sum_j q_jt^j\Bigr|
= \sum_{j \geq 1} q_j t^j
\leq \sum_{j\geq1}t^j = \frac{t}{1 - t} .
\]
于是
\[
\abs{p_0^{(n)} - p_0}
\leq \frac{2t}{1 - t}
+ \abs{G_{X_n}(t) - G_X(t)} .
\]
给定 $\varepsilon > 0$，取 $t$ 使 $\frac{2t}{1-t} <
\frac\varepsilon2$，再取 $n_0$，使末项 $<
\frac\varepsilon2$ 对 $n \geq n_0$ 成立，故 $p_0^{(n)} \to p_0$。

\emph{归纳步。}设 $p_j^{(n)} \to p_j$ 对 $j < k$ 成立。
考虑\emph{平移后的}函数
\[
g_n(t) = \frac{G_{X_n}(t) - p^{(n)}_0}{t}
= \sum_{j\geq0} p^{(n)}_{j+1}t^j,
\qquad
g(t) = \frac{G_X(t) - p_0}{t} ,
\]
这些是子概率序列的生成函数，
$(p^{(n)}_{j+1})_j$（总质量 $\leq 1$，这恰是
$k = 0$ 论证所用的全部条件）。对固定的 $t \in \intoo{0}{1}$，$g_n(t)
\to g(t)$ 由假设及情形 $k = 0$ 得出。应用 $k =
0$ 的论证于 $g_n$，得 $p_1^{(n)} \to p_1$；反复进行
平移 $k$ 次即得 $p_k^{(n)} \to p_k$，对每个 $k$ 成立。（这
是莱维连续性定理的离散、初等的实例，
其针对特征函数的一般形式，是本科
三年级的一个重要结论。）
\end{solution}

\begin{solution}{exo:b2:genfun:10}
逐点地，$\frac{1 + (-1)^X}{2}$ 等于 $1$（$X$ 为偶数时），
而 $0$ 对应奇数，故取期望（转移公式）得
\[
\P(X \text{ 偶数}) = \frac{1 + \E\bigl((-1)^X\bigr)}2 =
\frac{1 + G_X(-1)}2 .
\]
泊松分布：$\frac{1 + \eu^{-2\lambda}}2 \to \frac12$，当
$\lambda$ 增大时。二项分布：$\frac{1 + (1 - 2p)^n}2$。两种情形下
$G_X(-1) \to 0$ 说明 $X$ 的奇偶性变为公平
硬币：分布散布于许多整数上，遗忘了其
奇偶性。
\end{solution}

\begin{solution}{exo:b2:genfun:11}
$t + \dots + t^6 = t\,\frac{1 - t^6}{1 - t}$，且 $1 - t^6 =
(1 - t)(1 + t)(1 + t + t^2)(1 - t + t^2)$，由此得出所述的
因式分解。对第一颗骰子，$(1 + t)(1 + t + t^2) = 1 +
2t + 2t^2 + t^3$，故 $\frac{t(1+t)(1+t+t^2)}6 = \frac{t +
2t^2 + 2t^3 + t^4}6$：各面为 $\{1, 2, 2, 3, 3, 4\}$。对
第二颗骰子，展开
\[
(1 + 2t + 2t^2 + t^3)(1 - t + t^2)^2 = 1 + t^2 + t^3 + t^4 +
t^5 + t^7,
\]
故 $\frac{t(1+t)(1+t+t^2)(1-t+t^2)^2}6 = \frac{t + t^3 + t^4
+ t^5 + t^6 + t^8}6$：各面为 $\{1, 3, 4, 5, 6, 8\}$。
两个母函数之积将六个因子
重新组合为 $\bigl(\frac{t(1+t)(1+t+t^2)(1-t+t^2)}6
\bigr)^2$，即标准骰子母函数之平方：
西彻曼骰子对的总和恰有标准分布 ---
母函数可刻画所有此类重新分组。
\end{solution}

\begin{solution}{exo:b2:genfun:12}
设 $A$ 与 $B$ 分别为从
“当前无正面”与“当前有一个正面”出发的剩余时长的母函数。
掷一次耗去一步，于是：从状态 $0$ 出发，
反面回到状态 $0$，正面转入状态 $1$；从状态
$1$ 出发，正面结束游戏，反面回到状态 $0$：
\[
A(t) = t\bigl(q\,A(t) + p\,B(t)\bigr),
\qquad
B(t) = t\bigl(p + q\,A(t)\bigr).
\]
代入得 $A(1 - qt) = pt\,B = pt(pt + qtA)$，于是
\[
G_T(t) = A(t) = \frac{p^2t^2}{1 - qt - pq\,t^2} .
\]
当 $t = 1$ 时，分母为 $1 - q - pq = p(1 - q) = p^2$：
$G_T(1) = 1$，即游戏几乎必然结束（正如
\cref{exo:b2:proba:6}用递推所证）。在 $1$ 处取对数求导：
$\E(T) = 2 - \frac{D'(1)}{D(1)}$，其中
$D(t) = 1 - qt - pqt^2$，$D'(1) = -q - 2pq$：
\[
\E(T) = 2 + \frac{q + 2pq}{p^2} = \frac{2p^2 + q + 2pq}{p^2}
= \frac{1 + p}{p^2},
\]
它等于 $6$，此时 $p = \frac12$。
\end{solution}

\begin{solution}{pb:b2:genfun:1}
\textbf{1.} 对 $t \in \intoo01$，由 $G_n = G
\circ G_{n-1}$ 的链式法则得 $G_n'(t) =
G'\bigl(G_{n-1}(t)\bigr)G_{n-1}'(t)$。当 $t \to 1^-$ 时，
$G_{n-1}(t) \uparrow 1$，$G'$ 非减，
且其左极限为 $m$（在 $1$ 处），故第一个因子趋于 $m$。
用归纳法得第二个因子趋于 $m^{n-1}$。由
\cref{thm:b2:genfun:moments}可知 $\E(Z_n) = G_n'(1^-) = m^n$。

\textbf{2.} 再求一次导数，
\[
G_n'' = G''(G_{n-1})\,(G_{n-1}')^2 +
G'(G_{n-1})\,G_{n-1}'',
\]
令 $t \to 1^-$，得 $a_n = G''(1)m^{2(n-1)} + m\,
a_{n-1}$，其中 $a_n = G_n''(1)$，$a_1 = G''(1)$。对 $m \neq
1$ 时，由归纳可验证 $a_n = G''(1)\,m^{n-1}
\frac{m^n - 1}{m - 1}$（递推式即将 $G''(1)m^{2n-2}$
加到 $m\cdot G''(1)m^{n-2}\frac{m^{n-1}-1}{m-1}$上，且
$m^{n-1} + \frac{m^{n-1}-1}{m-1} = \frac{m^n - 1}{m-1}$）；
当 $m = 1$ 时，$a_n = a_{n-1} + G''(1) = n\,G''(1)$。

\textbf{3.} $V(Z_n) = a_n + m^n - m^{2n}$，且 $G''(1) =
\sigma^2 + m^2 - m$。当 $m \neq 1$ 时，项 $(m^2 -
m)m^{n-1}\frac{m^n-1}{m-1} = m^n(m^n - 1)$ 恰与 $m^n -
m^{2n}$ 相抵，剩下 $V(Z_n) =
\sigma^2m^{n-1}\frac{m^n-1}{m-1}$。当 $m = 1$ 时，$V(Z_n) =
nG''(1) = n\sigma^2$。

\textbf{4.} $Z_n$ 是非负整数值随机变量，故
$\P(Z_n > 0) = \P(Z_n \geq 1) \leq \E(Z_n) = m^n$，由马尔可夫
不等式（\cref{thm:b2:randomvar:markov}）。当 $m < 1$ 时，此概率
几何衰减——而且可和，故由博雷尔–坎泰利引理甚至可得
仅有有限多个世代非空，
这正是灭绝。

\textbf{5.} 柯西–施瓦茨不等式：$\E(Z_n)^2 = \E(Z_n\mathbf
1_{Z_n>0})^2 \leq \E(Z_n^2)\,\P(Z_n > 0)$。结合第 3 问
及 $m < 1$：
\[
\E(Z_n^2) = V(Z_n) + m^{2n}
\leq \frac{\sigma^2m^{n-1}}{1-m} + m^{2n},
\]
于是，将 $m^{2n}$ 除以该上界，并约去
$m^n$，得
\[
\P(Z_n > 0) \geq \frac{m^n}{\frac{\sigma^2}{m(1-m)} + m^n}
\geq \Bigl(\frac{\sigma^2}{m(1-m)} + 1\Bigr)^{-1}m^n ,
\]
利用分母中的 $m^n \leq 1$。结合第 4 问，
$\P(Z_n > 0) \asymp m^n$。

\textbf{6.} $G(t) = q\sum_k(pt)^k = \frac{q}{1 - pt}$，且 $m
= G'(1) = \frac{pq}{(1-p)^2} = \frac pq$。当 $p
< \frac12$ 时为次临界，$p = \frac12$ 时为临界，$p
> \frac12$ 时为超临界。

\textbf{7.} 方程 $G(t) = t$ 即 $pt^2 - t + q = 0$，其根为
$\frac{1 \pm \abs{p - q}}{2p}$，即 $1$ 与 $\frac qp =
\frac1m$。灭绝概率是
$\intcc01$ 中最小的不动点（\cref{thm:b2:genfun:extinction}）：
$q_{\mathrm{ext}} = 1$，若 $m \leq 1$；$\frac1m$，若 $m >
1$。

\textbf{8.} 对 $m \neq 1$，取 $p = \frac m{m+1}$，$q =
\frac1{m+1}$。若 $q_n = \frac{m^n - 1}{m^{n+1} - 1}$，则
\[
1 - p\,q_n = \frac{(m+1)(m^{n+1} - 1) - m(m^n - 1)}
{(m+1)(m^{n+1} - 1)} = \frac{m^{n+2} - 1}{(m+1)(m^{n+1} -
1)},
\]
于是 $q_{n+1} = \frac{q}{1 - pq_n} = \frac{m^{n+1} -
1}{m^{n+2} - 1}$，归纳基础 $q_0 = 0$ 成立。当 $m = 1$ 时：
$G(t) = \frac1{2 - t}$，且 $q_{n+1} = \frac1{2 -
\frac{n}{n+1}} = \frac{n+1}{n+2}$，$q_0 = 0$。

\textbf{9.} $1 - q_n = \frac{m^n(m - 1)}{m^{n+1} - 1}$。当
$m < 1$ 时，分母趋于 $-1$，故 $1 - q_n \sim (1 -
m)\,m^n$。当 $m > 1$ 时：
\[
q_{\mathrm{ext}} - q_n = \frac1m - \frac{m^n - 1}{m^{n+1} -
1} = \frac{m - 1}{m\,(m^{n+1} - 1)} \sim \frac{m -
1}{m^{2}}\;m^{-n} .
\]
又 $G'(t) = \frac{pq}{(1 - pt)^2}$ 在 $t = \frac
qp$ 处（此时 $1 - pt = 1 - q = p$）取值为 $G'(q_{\mathrm{ext}})
= \frac qp = \frac1m$，即所观察到的比值 $m^{-1}$ 恰为
吸引不动点处的导数。

\textbf{10.} 当 $p = \frac12$ 时：$G''(t) = \frac{1/4}{(1 -
t/2)^3}$，故 $G''(1) = 2$，且 $\sigma^2 = G''(1) + m - m^2 =
2$。由闭式得 $1 - q_n = \frac1{n+1}$：存活概率
按 $1/n$ 的速度衰减——过慢以致不可和，
不同于任何亚临界速率。

\textbf{11.} 归纳法：$G_1(t) = \frac1{2-t}$ 与
$n = 1$ 时的公式相符，且
\[
G(G_n(t)) = \cfrac{1}{2 - \cfrac{n - (n-1)t}{n+1 - nt}}
= \frac{n + 1 - nt}{2(n+1) - 2nt - n + (n-1)t}
= \frac{n+1 - nt}{n + 2 - (n+1)t} .
\]
于是
\[
\frac{G_n(t) - q_n}{1 - q_n}
= (n+1)\,\Bigl(\frac{n - (n-1)t}{n+1 - nt} -
\frac{n}{n+1}\Bigr)
= \frac{t}{n + 1 - nt}
= \frac{\frac{t}{n+1}}{1 - \frac{n}{n+1}t} ,
\]
即几何分布 $\mathcal G\bigl(\frac1{n+1}
\bigr)$ 在 $\N^*$ 上的概率母函数（见\cref{ex:b2:genfun:classical}）：已知存活，
$\P(Z_n = k \mid Z_n > 0) =
\frac1{n+1}\bigl(\frac n{n+1}\bigr)^{k-1}$，条件均值
为 $n + 1$。无条件均值 $1 = \E(Z_n)$ 恰为
存活概率趋于零与条件规模线性增长的
乘积。

\textbf{12.} 若该谱系在第 $n$ 代灭绝，则
$Y = Z_0 + \dots + Z_{n-1}$ 有限；若永不灭绝，
则 $Y \geq \sum_n 1 = \infty$。因此 $\{Y < \infty\}$ 为
灭绝事件，且 $\P(Y < \infty) = q_{\mathrm{ext}}$。
$H(t) = tG(H(t))$ 的推导在
\cref{exo:b2:genfun:7} 中——祖先贡献
因子 $t$，其子代各自找到 $Y$ 的独立副本，
经 $G$ 计数——仅用到
\cref{thm:b2:genfun:compound}，该式在各种情形下均成立。

\textbf{13.} 取 $G(s) = \frac{1 + s^2}2$，方程化为
$tH^2 - 2H + t = 0$，故 $H = \frac{1 - \sqrt{1 - t^2}}{t}$
（即满足 $H(0) = 0$ 的根）。与卡特兰级数
$C(x) = \frac{1 - \sqrt{1 - 4x}}{2x}$
（见\cref{ex:b2:powerseries:catalan}）比较，得 $H(t) = \frac
t2\,C\bigl(\frac{t^2}4\bigr) =
\sum_{k\geq0}C_k\,\frac{t^{2k+1}}{2^{2k+1}}$，即 $\P(Y =
2k+1) = C_k2^{-2k-1}$。检验：$\P(Y = 1) = C_0/2 = \frac12$
（祖先没有子代）；$\P(Y = 3) = C_1/8 = \frac18$
（两个子代，且均无子代：$\frac12\cdot\frac12\cdot
\frac12$）。

\textbf{14.} 对 $H = tG(H)$ 在 $\intoo01$ 上求导，并
令 $t \to 1^-$（取单调极限，见
\cref{thm:b2:genfun:moments}）：$H'(1)\bigl(1 - G'(H(1))\bigr)
= G(H(1))$。亚临界情形下 $H(1) = 1$，且 $\E(Y) =
H'(1) = \frac1{1 - m}$。临界情形下 $G'(1) = 1$
使左端因子为零，而右端为 $1$：无
有限的 $H'(1)$ 存在，故 $\E(Y) = \infty$——然而 $\P(Y <
\infty) = q = 1$。

\textbf{15.} $C_k = \frac1{k+1}\binom{2k}k \sim
\frac{4^k}{\sqrt\pi\,k^{3/2}}$，由
\cref{ex:b2:comparison:centralbinomial}，可得
\[
\P(Y = 2k+1) = \frac{C_k}{2\cdot4^{k}} \sim
\frac1{2\sqrt\pi\,k^{3/2}} .
\]
对尾部求和（与 $\int_K^\infty
k^{-3/2}\dd k = 2K^{-1/2}$，上、下两侧）：$\P(Y > 2K)
\asymp K^{-1/2}$，即 $\P(Y > n) \asymp n^{-1/2}$，是
一种期望为无穷的重尾，刻画了第 14 题中的问题。

\textbf{16.} 两个临界对象——公平随机游走的返回
时间（\cref{ch:b2:proba}中的周末问题）与
临界总后代数——几乎必然有限，
期望为无穷，局部律的指数为 $-3/2$，尾部
的指数为 $-1/2$。这并非偶然：探索一棵
家族树（逐个子代）产生一条 $\pm1$ 路径（每次出生上升一步，
每次死亡下降一步），恰为一个公平
随机游走，而 $Y$ 成为首次通过时间。临界性意味着
零漂移：过程始终处于
灭绝与爆炸的边缘，而 $\sqrt{}$ 尺度的
零漂移随机性的涨落恰好产生这些
指数。

\textbf{17.} $G''(t) = \sum_{n\geq2}n(n-1)p_nt^{n-2}$ 的各项非负，
故它在 $\intco01$ 上单调不减，并且
有限极限 $G''(1) = \sigma^2$（临界性使得
$\E Z_1(Z_1 - 1) = \sigma^2$）；单调不减函数
若极限等于边界值，则在 $1$ 处连续。
在点 $1$ 处用带积分余项的泰勒公式：
\[
G(t) = 1 + (t - 1) + \int_1^t(t - s)G''(s)\,\dd s
= t + \frac{G''(1)}2(1-t)^2 + o\bigl((1-t)^2\bigr),
\]
因 $G''(s) = G''(1) + o(1)$（$s \to 1^-$）。

\textbf{18.} 通分后，$h(t) =
\frac{G(t) - t}{(1 - G(t))(1 - t)}$。由第 17 题，
分子为 $b(1-t)^2 + o((1-t)^2)$，且 $1 - G(t) = (1 -
t)\bigl(1 - b(1-t) + o(1-t)\bigr)$，故 $h(t) \to b$。

\textbf{19.} 由 $h$ 在 $t = q_j$ 处的定义及 $G(q_j) =
q_{j+1}$：$\frac1{1 - q_{j+1}} - \frac1{1-q_j} = h(q_j)$；
自 $j = 0$（$q_0 = 0$）起求和，即得上式。由于
临界过程灭绝，$q_j \uparrow 1$，故 $h(q_j) \to
b$，且切萨罗均值 $\frac1n\sum_{j<n}h(q_j) \to b$：
$\frac1{1-q_n} \sim bn$，即
\[
\P(Z_n > 0) \sim \frac1{bn} = \frac{2}{\sigma^2 n} .
\]

\textbf{20.} 几何临界情形：$\sigma^2 = 2$
（第 10 问），故柯尔莫哥洛夫预言 $1 - q_n \sim \frac1n$
--- 而精确值为 $\frac1{n+1}$。

\textbf{21.} 由于 $Z_n\mathbf 1_{Z_n > 0} = Z_n$，$\E(Z_n
\mid Z_n > 0) = \frac{\E(Z_n)}{\P(Z_n > 0)} = \frac1{1 -
q_n} \sim \frac{\sigma^2n}2$。此时几何情形下的值为
$n + 1$，与第 11 问完全吻合（$\sigma^2 = 2$）。
临界图景：灭绝必然发生，平均规模
保持在 $1$，而稀有的存活谱系的规模
线性增长 --- 两者相互平衡。

\textbf{22.} 对于 $\mathcal P(\lambda)$ 后代，$G(t) =
\eu^{\lambda(t-1)}$，而灭绝概率是
方程 $q = \eu^{\lambda(q-1)}$ 的最小根。数值计算如下：
当 $\lambda = 1.5$ 时，$q \approx 0.417$（迭代 $q \mapsto
\eu^{1.5(q-1)}$：$0, 0.223, 0.312, 0.356, \dots \to 0.4172$）；
当 $\lambda = 2$ 时，$q \approx 0.203$。因此，一个首例病例引发大规模暴发的概率
为 $58\%$
（$\lambda = 1.5$），或 $80\%$（$\lambda = 2$）--- 很可能，但并非
必然。从 $q_0 = 0$ 出发的迭代收敛于
\emph{最小}根，因为 $G$ 单调不减。由归纳法
$q_n \leq r$（对任一不动点 $r$），且 $(q_n)$
递增（它是 $\P(Z_n = 0)$），故其极限是一个不动点，
且小于其他一切不动点。

\textbf{23.} 这 $k$ 个祖先各自开出独立的家族树，
而总灭绝是 $k$ 个独立灭绝事件的交，
概率为 $q^k$。对于
$\lambda = 1.5$：暴发概率 $1 - q^k \geq 0.99$
要求 $q^k \leq 0.01$，即 $k \geq
\frac{\ln 0.01}{\ln 0.417} \approx 5.3$：六个初始病例
使暴发的把握达 $99\%$。

\textbf{24.} \emph{$G'(q) < 1$：} $G - \mathrm{id}$ 是凸函数，
且在 $q$ 与 $1$ 处为零，故 $\leq 0$ 在 $\intcc
q1$ 上成立；若 $G'(q) = 1$，则 $q$ 处的切线（凸性使之
位于 $G$ 之下）将迫使 $G(t) \geq t$ 于 $\intcc q1$ 上，
从而 $G \equiv \mathrm{id}$ 于该区间，使所有系数
$p_n$（$n \geq 2$）为零，与 $m > 1$ 矛盾。\emph{几何
收敛：} $q_n < q$ 对一切 $n$ 成立（归纳法，$G$
递增），且由中值定理得 $q - q_{n+1} =
G'(c_n)(q - q_n)$，其中 $c_n \in \intoo{q_n}q$，故 $G'(c_n)
\leq G'(q) < 1$，且 $q - q_n \leq q\,G'(q)^n$。
\emph{伴随过程：} $\widehat G(t) = G(qt)/q =
\sum_kp_kq^{k-1}t^k$ 的系数非负，且
$\widehat G(1) = G(q)/q = 1$：概率母函数；其均值为 $\widehat
G'(1) = G'(q) < 1$，为次临界。几何族：$G(t) =
\frac{q}{1-pt}$，$q_{\mathrm{ext}} = \frac qp$，且
\[
\widehat G(t) = \frac pq\cdot\frac{q}{1 - p\frac qp t}
= \frac{p}{1 - qt} :
\]
即 $p$ 与 $q$ 互换后的几何后代分布 ---
在其灭绝事件上所见的超临界过程，正是
那个次临界的镜像过程。

\textbf{25.} 表格如下：$m < 1$：$q = 1$，$\P(Z_n > 0)
\asymp m^n$（第 4--5 问），$\E(Y) = \frac1{1-m}$，
存活世代的条件均值有界。$m = 1$：
$q = 1$，$\P(Z_n > 0) \sim \frac2{\sigma^2n}$（柯尔莫哥洛夫），
$\E(Y) = \infty$ 且 $\P(Y > n) \asymp n^{-1/2}$，存活者
规模为 $\sim \frac{\sigma^2n}2$。$m > 1$：$q < 1$ 是
最小不动点，$q - q_n = O(G'(q)^n)$，$\E(Z_n) = m^n$
增长，且在以灭绝为条件时，该过程即为
次临界伴随过程（第 24 问）。所用工具有：
概率母函数的复合把种群递推化为函数迭代；
凸性确定了不动点的几何结构；泰勒展开
在 $1^-$ 处将矩假设转化为局部展开；并且
切萨罗平均从一个裂项和中提取出柯尔莫哥洛夫的 $1/n$。
本科三年级分册补充了鞅 $Z_n/
m^n$ --- 其几乎必然极限把 $\E(Z_n) = m^n$ 细化为
逐条轨道的增长速率 --- 以及亚格洛姆
定理，即条件几何背后的极限分布，
它解释了第 11 问中所观察到的现象。
\end{solution}
